\subsection{MINIMAL CAUCHY FILTERS}

The minimal elements (with respect to inclusion) of the set of Cauchy filters on a uniform space X are called minimal Cauchy filters on X.

PROPOSITION 5. Let X be a uniform space. For each Cauchy filter F on X there is a unique minimal Cauchy filter \(\mathfrak{F}_0\) coarser than F. If B is a base of F and S is a fundamental system of symmetric entourages of X, then the sets V(M) (M ∈ B, V ∈ S) form a base of \(\mathfrak{F}_0\) .

If M, \(M'\) are in B and V, \(V'\) are in S, then there is a set \(M'' \in B\) (resp. \(V'' \in S\) ) such that \(M'' \subset M \cap M'\) (resp. \(V'' \subset V \cap V'\) ); hence \(V''(M'') \subset V(M) \cap V'(M')\) and therefore the sets \(V(M)\) ( \(M \in B\) , \(V \in S\) ) indeed form a base of a filter \(F_0\) on X. Further, if M is V-small,

then V(M) is V-small; hence \(\mathfrak{F}_0\) is a Cauchy filter and is clearly coarser than \(\mathfrak{F}\) . To complete the proof it is enough to show that if \(\mathfrak{G}\) is a Cauchy filter coarser than \(\mathfrak{F}\) , then \(\mathfrak{G}\) is finer than \(\mathfrak{F}_0\) . For each M ∈ \(\mathfrak{B}\) and each V ∈ \(\mathfrak{S}\) there is a set N ∈ \(\mathfrak{G}\) which is V-small; since N ∈ \(\mathfrak{F}\) , N meets M; hence N ⊂ V(M) and so V(M) ∈ \(\mathfrak{G}\) .

COROLLARY 1. For each \(x \in \mathbf{X}\) , the neighbourhood filter \(\mathfrak{B}(x)\) of \(x\) in \(\mathbf{X}\) is a minimal Cauchy filter.

Take \(\mathfrak{F}\) in Proposition 5 to be the filter of all subsets of \(\mathbf{X}\) which contain \(x\) , and take \(\mathfrak{B}\) to consist of the single element \(\{x\}\) .

COROLLARY 2. Every cluster point x of a Cauchy filter F is a limit point of F.

There is a filter \(\mathfrak{G}\) which is finer than both \(\mathfrak{F}\) and \(\mathfrak{B}(x)\) (Chapter I, § 7, no. 2, Proposition 4); since \(\mathfrak{F}\) is a Cauchy filter, so is \(\mathfrak{G}\) . If \(\mathfrak{F}_0\) is the unique minimal Cauchy filter coarser than \(\mathfrak{F}\) , then both \(\mathfrak{F}_0\) and \(\mathfrak{B}(x)\) are minimal Cauchy filters coarser than \(\mathfrak{G}\) . Hence \(\mathfrak{F}_0 = \mathfrak{B}(x)\) , which shows that \(\mathfrak{F}\) converges to \(x\) .

COROLLARY 3. Every Cauchy filter, which is coarser than a filter converging to a point x, also converges to x.

This is a consequence of Corollary 2.

COROLLARY 4. If \(\mathfrak{F}\) is a minimal Cauchy filter, then every set of \(\mathfrak{F}\) has a non-empty interior which also belongs to \(\mathfrak{F}\) (in other words, \(\mathfrak{F}\) has a base consisting of open sets).

Let V be any entourage of X; then there is an open entourage \(U \subset V\) (§ 1, no. 2, Corollary 2 to Proposition 2). For each subset M of X, \(U(M)\) is open and contained in \(V(M)\) ; hence the result, in view of Proposition 5.

